\subsection{\label{sec.sign.pie}The principles of inclusion and exclusion}

We have so far been applying the \textquotedblleft yoga\textquotedblright\ of
sign-reversing involutions directly to alternating sums. However, some of its
essence can also be crystallized into useful theorems. Most famous among such
theorems are the \emph{principles of inclusion and exclusion} (also known as
\emph{Sylvester sieve theorems} or \emph{Poincar\'{e}'s theorems}%
)\footnote{People typically use the singular forms (\textquotedblleft
principle\textquotedblright\ and \textquotedblleft theorem\textquotedblright%
\ rather than \textquotedblleft principles\textquotedblright\ and
\textquotedblleft theorems\textquotedblright), but often mean different
things.}.

\subsubsection{\label{subsec.sign.pie.size}The size version}

The simplest such principle answers the following question: Assume that you
have a finite set $U$, and some subsets $A_{1},A_{2},\ldots,A_{n}$ of $U$.
Assume that, for each selection of some of these subsets, you know how many
elements of $U$ belong to all selected sets at the same time. (For instance,
you know how many elements of $U$ belong to $A_{2}$, $A_{3}$ and $A_{5}$ at
the same time.) Does this help you count the elements of $U$ that belong to
\textbf{none} of the $n$ subsets (i.e., that don't belong to $A_{1}\cup
A_{2}\cup\cdots\cup A_{n}$) ?

The answer is \textquotedblleft yes\textquotedblright, and there is an
explicit formula for this count:

\begin{theorem}
[size version of the PIE]\label{thm.pie.1}Let $n\in\mathbb{N}$. Let $U$ be a
finite set. Let $A_{1},A_{2},\ldots,A_{n}$ be $n$ subsets of $U$. Then,%
\begin{align*}
&  \left(  \text{\# of }u\in U\text{ that satisfy }u\notin A_{i}\text{ for all
}i\in\left[  n\right]  \right) \\
&  =\sum_{I\subseteq\left[  n\right]  }\left(  -1\right)  ^{\left\vert
I\right\vert }\left(  \text{\# of }u\in U\text{ that satisfy }u\in A_{i}\text{
for all }i\in I\right)  .
\end{align*}

\end{theorem}

Some explanations are in order:

\begin{itemize}
\item Here and in the following, we are using the notation $\left[  n\right]
$ for the set $\left\{  1,2,\ldots,n\right\}  $, as defined in Definition
\ref{def.perm.Sn-iven}.

\item The summation sign \textquotedblleft$\sum_{I\subseteq\left[  n\right]
}$\textquotedblright\ means a sum over all subsets $I$ of $\left[  n\right]
$. More generally, if $S$ is a given set, then the summation sign
\textquotedblleft$\sum_{I\subseteq S}$\textquotedblright\ shall always mean a
sum over all subsets $I$ of $S$.

\item The shorthand \textquotedblleft\emph{PIE}\textquotedblright\ in the name
of Theorem \ref{thm.pie.1} is short for \textquotedblleft\emph{Principle of
Inclusion and Exclusion}\textquotedblright.
\end{itemize}

In one form or another, Theorem \ref{thm.pie.1} appears in almost any text on
combinatorics (e.g., in \cite[Theorem 2.1.1]{Sagan19}, in \cite[\S 4.11]%
{Loehr-BC}, in \cite[Theorem 5.3]{Strick20}, in \cite[Theorem 2.9.7]{19fco},
or -- in almost the same form as above -- in \cite[Theorem 7.8.6]{20f}). Most
commonly, its claim is stated in the shorter (if less transparent) form%
\[
\left\vert U\setminus\left(  A_{1}\cup A_{2}\cup\cdots\cup A_{n}\right)
\right\vert =\sum_{I\subseteq\left[  n\right]  }\left(  -1\right)
^{\left\vert I\right\vert }\left\vert \bigcap\limits_{i\in I}A_{i}\right\vert
,
\]
where $\bigcap\limits_{i\in I}A_{i}$ denotes the set $\left\{  u\in
U\ \mid\ u\in A_{i}\text{ for all }i\in I\right\}  $ whenever $I$ is a subset
of $\left[  n\right]  $\ \ \ \ \footnote{This notation should be taken with a
grain of salt. When $I$ is a \textbf{nonempty} subset of $\left[  n\right]  $,
it is indeed true that the set $\bigcap\limits_{i\in I}A_{i}$ as we just
defined it (i.e., the set $\left\{  u\in U\ \mid\ u\in A_{i}\text{ for all
}i\in I\right\}  $) is the intersection of the sets $A_{i}$ over all $i\in I$
(that is, if $I=\left\{  i_{1},i_{2},\ldots,i_{k}\right\}  $, then
$\bigcap\limits_{i\in I}A_{i}=A_{i_{1}}\cap A_{i_{2}}\cap\cdots\cap A_{i_{k}}%
$). However, if $I$ is the \textbf{empty} set, then the literal intersection
of the sets $A_{i}$ over all $i\in I$ is not well-defined (indeed, by common
sense, such an intersection should contain every object whatsoever; but there
is no set that does this), whereas the set we just called $\bigcap
\limits_{i\in I}A_{i}$ is simply the set $U$. This is still justified, since
we should think of the sets $A_{i}$ not as arbitrary sets but rather as
subsets of $U$ (so that any intersections are to be taken within $U$). This,
incidentally, is the reason for the choice of the letter \textquotedblleft%
$U$\textquotedblright: We think of $U$ as the \textquotedblleft
universe\textquotedblright\ in which our objects live.}. This form is indeed
equivalent to the claim of Theorem \ref{thm.pie.1}, since we have%
\begin{align*}
&  \left\vert U\setminus\left(  A_{1}\cup A_{2}\cup\cdots\cup A_{n}\right)
\right\vert \\
&  =\left(  \text{\# of }u\in U\text{ that satisfy }u\notin A_{1}\cup
A_{2}\cup\cdots\cup A_{n}\right) \\
&  =\left(  \text{\# of }u\in U\text{ that satisfy }u\notin A_{i}\text{ for
all }i\in\left[  n\right]  \right)
\end{align*}
and since (for each subset $I$ of $\left[  n\right]  $) we have%
\[
\left\vert \bigcap\limits_{i\in I}A_{i}\right\vert =\left(  \text{\# of }u\in
U\text{ that satisfy }u\in A_{i}\text{ for all }i\in I\right)
\]
(by the definition of $\bigcap\limits_{i\in I}A_{i}$ that we just gave).

Rather than prove Theorem \ref{thm.pie.1} directly, we shall soon derive it
from more general results (in order to avoid duplicating arguments). First,
however, let us give an interpretation that makes Theorem \ref{thm.pie.1} a
little bit more intuitive, and sketch four applications (more can be found in
textbooks -- e.g., \cite[\S 2.1]{Sagan19}, \cite[Chapter 2]{Stanley-EC1},
\cite[Chapter 3]{Wildon19}, \cite[Chapter 6]{AndFen04}, \cite[\S 2.9]{19fco}).

\begin{statement}
\textbf{\textquotedblleft Rule-breaking\textquotedblright\ interpretation of
Theorem \ref{thm.pie.1}.} Assume that we are given a finite set $U$, and we
are given $n$ rules (labelled $1,2,\ldots,n$) that each element of $U$ may or
may not satisfy. (For instance, one element of $U$ might satisfy all of these
rules; another might satisfy none; yet another might satisfy rules $1$ and $3$
only. A rule can be something like \textquotedblleft thou shalt be divisible
by $5$\textquotedblright\ (if the elements of $U$ are numbers) or
\textquotedblleft thou shalt be a nonempty set\textquotedblright\ (if the
elements of $U$ are sets).)

Assume that, for each $I\subseteq\left[  n\right]  $, we know how many
elements $u\in U$ satisfy all rules in $I$ (but may or may not satisfy the
remaining rules). For example, this means that we know how many elements $u\in
U$ satisfy rules $2,3,5$ (simultaneously). Then, we can compute the \# of
elements $u\in U$ that violate all $n$ rules $1,2,\ldots,n$ by the following
formula:%
\begin{align*}
&  \left(  \text{\# of elements }u\in U\text{ that violate all }n\text{ rules
}1,2,\ldots,n\right) \\
&  =\sum_{I\subseteq\left[  n\right]  }\left(  -1\right)  ^{\left\vert
I\right\vert }\left(  \text{\# of elements }u\in U\text{ that satisfy all
rules in }I\right)  .
\end{align*}


Indeed, this formula is precisely what we obtain if we apply Theorem
\ref{thm.pie.1} to the $n$ subsets $A_{1},A_{2},\ldots,A_{n}$ defined by
setting%
\[
A_{i}:=\left\{  u\in U\ \mid\ u\text{ satisfies rule }i\right\}
\ \ \ \ \ \ \ \ \ \ \text{for each }i\in\left[  n\right]  .
\]

\end{statement}

Thus, if you have a counting problem that can be restated as \textquotedblleft
count things that violate a bunch of rules\textquotedblright, then you can
apply Theorem \ref{thm.pie.1} (in the interpretation we just gave) to
\textquotedblleft turn the problem positive\textquotedblright, i.e., to make
it about counting rule-followers instead of rule-violators. If the
\textquotedblleft positive\textquotedblright\ problem is easier, then this is
a useful technique. We will now witness this on four examples.

\subsubsection{\label{subsec.sign.pie.exas}Examples}

\textbf{Example 1.} Let $n,m\in\mathbb{N}$. Let us compute the \# of
surjective maps from $\left[  m\right]  $ to $\left[  n\right]  $. (We will
outline the argument here; details can be found in \cite[\S 7.8.2]{20f} or
\cite[\S 2.9.4]{19fco}.)

What are surjective maps? They are maps that take each element of the target
set as a value. Thus, in particular, a map $f:\left[  m\right]  \rightarrow
\left[  n\right]  $ is surjective if and only if it takes each $i\in\left[
n\right]  $ as a value.

Hence, if we impose $n$ rules $1,2,\ldots,n$ on a map $f:\left[  m\right]
\rightarrow\left[  n\right]  $, where rule $i$ says \textquotedblleft thou
shalt not take $i$ as a value\textquotedblright, then the surjective maps
$f:\left[  m\right]  \rightarrow\left[  n\right]  $ are precisely the maps
$f:\left[  m\right]  \rightarrow\left[  n\right]  $ that violate all $n$
rules. Hence,%
\begin{align*}
&  \left(  \text{\# of surjective maps }f:\left[  m\right]  \rightarrow\left[
n\right]  \right) \\
&  =\left(  \text{\# of maps }f:\left[  m\right]  \rightarrow\left[  n\right]
\text{ that violate all }n\text{ rules }1,2,\ldots,n\right) \\
&  =\sum_{I\subseteq\left[  n\right]  }\left(  -1\right)  ^{\left\vert
I\right\vert }\left(  \text{\# of maps }f:\left[  m\right]  \rightarrow\left[
n\right]  \text{ that satisfy all rules in }I\right)
\end{align*}
(by the \textquotedblleft rule-breaking\textquotedblright\ interpretation of
Theorem \ref{thm.pie.1}).

Now, fix some subset $I$ of $\left[  n\right]  $. What is the \# of maps
$f:\left[  m\right]  \rightarrow\left[  n\right]  $ that satisfy all rules in
$I$ ? A map $f:\left[  m\right]  \rightarrow\left[  n\right]  $ satisfies all
rules in $I$ if and only if it takes none of the $i\in I$ as a value, i.e., if
all its values belong to $\left[  n\right]  \setminus I$. Thus, the maps
$f:\left[  m\right]  \rightarrow\left[  n\right]  $ that satisfy all rules in
$I$ are nothing but the maps from $\left[  m\right]  $ to $\left[  n\right]
\setminus I$. The \# of such maps is therefore $\left\vert \left[  n\right]
\setminus I\right\vert ^{\left\vert \left[  m\right]  \right\vert }=\left(
n-\left\vert I\right\vert \right)  ^{m}$ (since $I\subseteq\left[  n\right]  $
entails $\left\vert \left[  n\right]  \setminus I\right\vert =\left\vert
\left[  n\right]  \right\vert -\left\vert I\right\vert =n-\left\vert
I\right\vert $, and since $\left\vert \left[  m\right]  \right\vert =m$).

Forget that we fixed $I$. We thus have shown that for each subset $I$ of
$\left[  n\right]  $, we have%
\begin{align}
&  \left(  \text{\# of maps }f:\left[  m\right]  \rightarrow\left[  n\right]
\text{ that satisfy all rules in }I\right) \nonumber\\
&  =\left(  n-\left\vert I\right\vert \right)  ^{m}. \label{eq.exa.pie1.1.3}%
\end{align}
Substituting this into the above computation, we find%
\begin{align*}
&  \left(  \text{\# of surjective maps }f:\left[  m\right]  \rightarrow\left[
n\right]  \right) \\
&  =\sum_{I\subseteq\left[  n\right]  }\left(  -1\right)  ^{\left\vert
I\right\vert }\underbrace{\left(  \text{\# of maps }f:\left[  m\right]
\rightarrow\left[  n\right]  \text{ that satisfy all rules in }I\right)
}_{\substack{=\left(  n-\left\vert I\right\vert \right)  ^{m}\\\text{(by
(\ref{eq.exa.pie1.1.3}))}}}\\
&  =\sum_{I\subseteq\left[  n\right]  }\left(  -1\right)  ^{\left\vert
I\right\vert }\left(  n-\left\vert I\right\vert \right)  ^{m}=\sum_{k=0}%
^{n}\ \ \sum_{\substack{I\subseteq\left[  n\right]  ;\\\left\vert I\right\vert
=k}}\underbrace{\left(  -1\right)  ^{\left\vert I\right\vert }\left(
n-\left\vert I\right\vert \right)  ^{m}}_{\substack{=\left(  -1\right)
^{k}\left(  n-k\right)  ^{m}\\\text{(since }\left\vert I\right\vert
=k\text{)}}}\\
&  \ \ \ \ \ \ \ \ \ \ \ \ \ \ \ \ \ \ \ \ \left(
\begin{array}
[c]{c}%
\text{here, we have split the sum according to}\\
\text{the value of }\left\vert I\right\vert
\end{array}
\right) \\
&  =\sum_{k=0}^{n}\ \ \underbrace{\sum_{\substack{I\subseteq\left[  n\right]
;\\\left\vert I\right\vert =k}}\left(  -1\right)  ^{k}\left(  n-k\right)
^{m}}_{\substack{=\dbinom{n}{k}\left(  -1\right)  ^{k}\left(  n-k\right)
^{m}\\\text{(since this is a sum of }\dbinom{n}{k}\\\text{many equal
addends)}}}=\sum_{k=0}^{n}\dbinom{n}{k}\left(  -1\right)  ^{k}\left(
n-k\right)  ^{m}\\
&  =\sum_{k=0}^{n}\left(  -1\right)  ^{k}\dbinom{n}{k}\left(  n-k\right)
^{m}.
\end{align*}


Thus, we have proved the following theorem:

\begin{theorem}
\label{thm.pie.count-sur}Let $n,m\in\mathbb{N}$. Then,%
\[
\left(  \text{\# of surjective maps }f:\left[  m\right]  \rightarrow\left[
n\right]  \right)  =\sum\limits_{k=0}^{n}\left(  -1\right)  ^{k}\dbinom{n}%
{k}\left(  n-k\right)  ^{m}.
\]

\end{theorem}

This is the simplest expression for this number. It has no product formula
(unlike the \# of injective maps $f:\left[  m\right]  \rightarrow\left[
n\right]  $, which is $n\left(  n-1\right)  \left(  n-2\right)  \cdots\left(
n-m+1\right)  $).

Before we move on to the next example, let us draw a few consequences from
Theorem \ref{thm.pie.count-sur}:

\begin{corollary}
\label{cor.pie.count-sur.cors}Let $n\in\mathbb{N}$. Then: \medskip

\textbf{(a)} We have $\sum\limits_{k=0}^{n}\left(  -1\right)  ^{k}\dbinom
{n}{k}\left(  n-k\right)  ^{m}=0$ for any $m\in\mathbb{N}$ satisfying $m<n$.
\medskip

\textbf{(b)} We have $\sum\limits_{k=0}^{n}\left(  -1\right)  ^{k}\dbinom
{n}{k}\left(  n-k\right)  ^{n}=n!$. \medskip

\textbf{(c)} We have $\sum\limits_{k=0}^{n}\left(  -1\right)  ^{k}\dbinom
{n}{k}\left(  n-k\right)  ^{m}\geq0$ for each $m\in\mathbb{N}$. \medskip

\textbf{(d)} We have $n!\mid\sum\limits_{k=0}^{n}\left(  -1\right)
^{k}\dbinom{n}{k}\left(  n-k\right)  ^{m}$ for each $m\in\mathbb{N}$.
\end{corollary}

\begin{proof}
[Proof of Corollary \ref{cor.pie.count-sur.cors} (sketched).]\textbf{(a)} Let
$m\in\mathbb{N}$ satisfy $m<n$. Theorem \ref{thm.pie.count-sur} shows that the
sum $\sum\limits_{k=0}^{n}\left(  -1\right)  ^{k}\dbinom{n}{k}\left(
n-k\right)  ^{m}$ equals the \# of surjective maps $f:\left[  m\right]
\rightarrow\left[  n\right]  $. However, there are no such maps (by the
Pigeonhole Principle for Surjections)\footnote{Indeed, the Pigeonhole
Principle for Surjections shows that there are no surjective maps from a
smaller set to a larger set. Thus, in particular, there are no surjective maps
from $\left[  m\right]  $ to $\left[  n\right]  $ (since $m<n$).}; hence, this
\# is $0$. Therefore, the sum is $0$. In other words, $\sum\limits_{k=0}%
^{n}\left(  -1\right)  ^{k}\dbinom{n}{k}\left(  n-k\right)  ^{m}=0$. This
proves Corollary \ref{cor.pie.count-sur.cors} \textbf{(a)}. \medskip

\textbf{(b)} Theorem \ref{thm.pie.count-sur} (applied to $m=n$) shows that the
sum $\sum\limits_{k=0}^{n}\left(  -1\right)  ^{k}\dbinom{n}{k}\left(
n-k\right)  ^{n}$ equals the \# of surjective maps $f:\left[  n\right]
\rightarrow\left[  n\right]  $. However, these maps are precisely the
permutations of $\left[  n\right]  $ (by the Pigeonhole Principle for
Surjections)\footnote{Indeed, the Pigeonhole Principle for Surjections shows
that any surjective map between two finite sets of the same size is bijective.
Thus, any surjective map $f:\left[  n\right]  \rightarrow\left[  n\right]  $
is bijective, and hence is a permutation of $\left[  n\right]  $. The converse
is true as well (i.e., any permutation of $\left[  n\right]  $ is a surjective
map $f:\left[  n\right]  \rightarrow\left[  n\right]  $). Thus, the surjective
maps $f:\left[  n\right]  \rightarrow\left[  n\right]  $ are precisely the
permutations of $\left[  n\right]  $.}; therefore, this \# is $n!$. Therefore,
this sum is $n!$. This proves Corollary \ref{cor.pie.count-sur.cors}
\textbf{(b)}. \medskip

\textbf{(c)} Let $m\in\mathbb{N}$. Theorem \ref{thm.pie.count-sur} shows that
the sum $\sum\limits_{k=0}^{n}\left(  -1\right)  ^{k}\dbinom{n}{k}\left(
n-k\right)  ^{m}$ equals the \# of surjective maps $f:\left[  m\right]
\rightarrow\left[  n\right]  $. Hence, this sum is $\geq0$. This proves
Corollary \ref{cor.pie.count-sur.cors} \textbf{(c)}. \medskip

\textbf{(d)} Let $m\in\mathbb{N}$. Theorem \ref{thm.pie.count-sur} shows that
the sum $\sum\limits_{k=0}^{n}\left(  -1\right)  ^{k}\dbinom{n}{k}\left(
n-k\right)  ^{m}$ equals the \# of surjective maps $f:\left[  m\right]
\rightarrow\left[  n\right]  $. Let
\[
U:=\left\{  \text{surjective maps }f:\left[  m\right]  \rightarrow\left[
n\right]  \right\}  .
\]
Thus, this sum equals $\left\vert U\right\vert $. Hence, it remains to show
that $n!\mid\left\vert U\right\vert $.

The argument that follows will use the language of group actions (although it
could be restated combinatorially)\footnote{For a refresher on group actions,
see (e.g.) \cite[\S 3.2]{Quinlan21} or \cite[\S 9.11--\S 9.15]{Loehr-BC} or
\cite[\S 6.7--\S 6.9]{Artin} or \cite[Fall 2018, Weeks 8--10]%
{Armstrong-561fa18sp19} or \cite[\S 6.1]{Aigner07}. When $G$ is a group, we
use the word \textquotedblleft$G$\emph{-set}\textquotedblright\ to mean a set
on which $G$ acts.}. If $\sigma\in S_{n}$ is any permutation and $f:\left[
m\right]  \rightarrow\left[  n\right]  $ is a surjective map, then the
composition $\sigma\circ f:\left[  m\right]  \rightarrow\left[  n\right]  $ is
again a surjective map. In other words, any $\sigma\in S_{n}$ and $f\in U$
satisfy $\sigma\circ f\in U$. Hence, the symmetric group $S_{n}$ acts on the
set $U$ by the rule\footnote{This is called \textquotedblleft acting by
post-composition\textquotedblright\ (since $\sigma\circ f$ is obtained from
$f$ by composing with $\sigma$ \textquotedblleft after\textquotedblright%
\ $f$).}%
\[
\sigma\cdot f=\sigma\circ f\ \ \ \ \ \ \ \ \ \ \text{for all }\sigma\in
S_{n}\text{ and }f\in U.
\]
(This is a well-defined group action, since composition of maps is associative.)

This turns $U$ into an $S_{n}$-set. It is not hard to see that this action of
$S_{n}$ on $U$ is free -- i.e., the stabilizer of each $f\in U$ is the trivial
subgroup $\left\{  \operatorname*{id}\right\}  $ of $S_{n}$%
\ \ \ \ \footnote{\textit{Proof.} Let $f\in U$. We must prove that the
stabilizer of $f$ is $\left\{  \operatorname*{id}\right\}  $.
\par
Let $\sigma$ belong to the stabilizer of $f$. Thus, $\sigma\circ f=f$. Now,
let $j\in\left[  n\right]  $. Recall that $f\in U$; hence, $f$ is surjective.
Thus, there exists some $i\in\left[  m\right]  $ such that $j=f\left(
i\right)  $. Consider this $i$. Now, from $j=f\left(  i\right)  $, we obtain
$\sigma\left(  j\right)  =\sigma\left(  f\left(  i\right)  \right)
=\underbrace{\left(  \sigma\circ f\right)  }_{=f}\left(  i\right)  =f\left(
i\right)  =j$. Now, forget that we fixed $j$. We thus have shown that
$\sigma\left(  j\right)  =j$ for each $j\in\left[  n\right]  $. In other
words, $\sigma=\operatorname*{id}$.
\par
Forget that we fixed $\sigma$. We thus have shown that any $\sigma$ in the
stabilizer of $f$ satisfies $\sigma=\operatorname*{id}$. In other words, the
stabilizer of $f$ is a subset of $\left\{  \operatorname*{id}\right\}  $.
Therefore, the stabilizer of $f$ must be $\left\{  \operatorname*{id}\right\}
$ (since $\operatorname*{id}$ is clearly in the stabilizer of $f$).}. Thus,
the Orbit-Stabilizer Theorem (see, e.g., \cite[Theorem 3.2.5]{Quinlan21})
shows that each orbit of this action has size%
\[
\left[  S_{n}:\left\{  \operatorname*{id}\right\}  \right]
=\underbrace{\left\vert S_{n}\right\vert }_{=n!}/\underbrace{\left\vert
\left\{  \operatorname*{id}\right\}  \right\vert }_{=1}=n!/1=n!.
\]
However, the orbits of this action form a partition of $U$ (since each element
of $U$ belongs to exactly one orbit). Thus, we have%
\[
\left\vert U\right\vert =\left(  \text{the sum of the sizes of the
orbits}\right)  =\left(  \text{\# of orbits}\right)  \cdot n!
\]
(since each orbit has size $n!$). This entails $n!\mid\left\vert U\right\vert
$, which is exactly what we wanted to show. This proves Corollary
\ref{cor.pie.count-sur.cors} \textbf{(d)}.
\end{proof}

We note that parts \textbf{(a)} and \textbf{(b)} of Corollary
\ref{cor.pie.count-sur.cors} can also be proved algebraically (see, e.g.,
\cite[Exercise 5.4.2 \textbf{(d)}]{20f} for an algebraic generalization of
Corollary \ref{cor.pie.count-sur.cors} \textbf{(a)}); but this is harder for
Corollary \ref{cor.pie.count-sur.cors} \textbf{(d)} and (to my knowledge)
impossible for Corollary \ref{cor.pie.count-sur.cors} \textbf{(c)}.

\bigskip

\textbf{Example 2.} (See \cite[\S 2.9.5]{19fco} for details.) We will use the
following definition:

\begin{definition}
\label{def.pie.dera}A \emph{derangement} of a set $X$ means a permutation of
$X$ that has no fixed points.
\end{definition}

Now, let $n\in\mathbb{N}$. How many derangements does $\left[  n\right]  $ have?

Before answering this question, we establish notations and a few examples.

Let $D_{n}$ be the \# of derangements of $\left[  n\right]  $. For example:

\begin{itemize}
\item The identity permutation $\operatorname*{id}\in S_{0}$ is a derangement,
since it has no fixed points (since $\left[  0\right]  =\varnothing$ has no
elements to begin with). Thus, $D_{0}=1$.

\item The identity permutation $\operatorname*{id}\in S_{1}$ is \textbf{not} a
derangement, since $\operatorname*{id}\left(  1\right)  =1$. Thus, $D_{1}=0$.

\item In the symmetric group $S_{2}$, the identity is \textbf{not} a
derangement, but the transposition $s_{1}=t_{1,2}$ is one. Thus, $D_{2}=1$.

\item In the symmetric group $S_{3}$, the derangements are the $3$-cycles
$\operatorname*{cyc}\nolimits_{1,2,3}$ and $\operatorname*{cyc}%
\nolimits_{1,3,2}$. Thus, $D_{3}=2$.
\end{itemize}

Here is a table of early values of $D_{n}$:%
\[%
\begin{tabular}
[c]{|c||c|c|c|c|c|c|c|c|c|c|c|}\hline
$n$ & $0$ & $1$ & $2$ & $3$ & $4$ & $5$ & $6$ & $7$ & $8$ & $9$ & $10$\\\hline
$D_{n}$ & $1$ & $0$ & $1$ & $2$ & $9$ & $44$ & $265$ & $1854$ & $14\ 833$ &
$133\ 496$ & $1\ 334\ 961$\\\hline
\end{tabular}
\ \ \ \ \ \ \ .
\]


\textbf{Note:} The number $D_{n}$ is also called the \emph{subfactorial} of
$n$ and is sometimes denoted by $!n$. Be careful with that notation: what is
$2!2$ ?

\bigskip

Let us now try to compute $D_{n}$ in general. Fix $n\in\mathbb{N}$, and set
$U:=S_{n}=\left\{  \text{permutations of }\left[  n\right]  \right\}  $. We
impose $n$ rules $1,2,\ldots,n$ on a permutation $\sigma\in U$, with rule $i$
being \textquotedblleft thou shalt leave the element $i$
fixed\textquotedblright\ (in other words, rule $i$ requires a permutation
$\sigma\in U$ to satisfy $\sigma\left(  i\right)  =i$). Now,%
\begin{align*}
D_{n}  &  =\left(  \text{\# of derangements of }\left[  n\right]  \right) \\
&  =\left(  \text{\# of permutations }u\in U\text{ that violate all }n\text{
rules }1,2,\ldots,n\right) \\
&  =\sum_{I\subseteq\left[  n\right]  }\left(  -1\right)  ^{\left\vert
I\right\vert }\left(  \text{\# of permutations }u\in U\text{ that satisfy all
rules in }I\right)
\end{align*}
(by the \textquotedblleft rule-breaking\textquotedblright\ interpretation of
Theorem \ref{thm.pie.1}).

Now, let $I$ be a subset of $\left[  n\right]  $. What is the \# of
permutations $u\in U$ that satisfy all rules in $I$ ? These permutations are
the permutations of $\left[  n\right]  $ that leave each $i\in I$ fixed (but
can do whatever they want with the remaining elements of $\left[  n\right]
$). Clearly, there are $\left(  n-\left\vert I\right\vert \right)  !$ such
permutations, since they are essentially the permutations of the $\left(
n-\left\vert I\right\vert \right)  $-element set $\left[  n\right]  \setminus
I$ (with the elements of $I$ tacked on as fixed points). (See \cite[Corollary
2.9.16]{19fco} for the technical details of this intuitively obvious argument.)

Forget that we fixed $I$. We thus have shown that%
\begin{align}
&  \left(  \text{\# of permutations }u\in U\text{ that satisfy all rules in
}I\right) \nonumber\\
&  =\left(  n-\left\vert I\right\vert \right)  ! \label{eq.exa.pie1.2.3}%
\end{align}
for any subset $I$ of $\left[  n\right]  $.

Thus, the above computation becomes%
\begin{align*}
D_{n}  &  =\sum_{I\subseteq\left[  n\right]  }\left(  -1\right)  ^{\left\vert
I\right\vert }\underbrace{\left(  \text{\# of permutations }u\in U\text{ that
satisfy all rules in }I\right)  }_{\substack{=\left(  n-\left\vert
I\right\vert \right)  !\\\text{(by (\ref{eq.exa.pie1.2.3}))}}}\\
&  =\underbrace{\sum_{I\subseteq\left[  n\right]  }}_{=\sum_{k=0}^{n}%
\ \ \sum_{\substack{I\subseteq\left[  n\right]  ;\\\left\vert I\right\vert
=k}}}\left(  -1\right)  ^{\left\vert I\right\vert }\left(  n-\left\vert
I\right\vert \right)  !=\sum_{k=0}^{n}\ \ \sum_{\substack{I\subseteq\left[
n\right]  ;\\\left\vert I\right\vert =k}}\underbrace{\left(  -1\right)
^{\left\vert I\right\vert }\left(  n-\left\vert I\right\vert \right)
!}_{\substack{=\left(  -1\right)  ^{k}\left(  n-k\right)  !\\\text{(since
}\left\vert I\right\vert =k\text{)}}}\\
&  =\sum_{k=0}^{n}\ \ \underbrace{\sum_{\substack{I\subseteq\left[  n\right]
;\\\left\vert I\right\vert =k}}\left(  -1\right)  ^{k}\left(  n-k\right)
!}_{=\dbinom{n}{k}\left(  -1\right)  ^{k}\left(  n-k\right)  !}=\sum_{k=0}%
^{n}\dbinom{n}{k}\left(  -1\right)  ^{k}\left(  n-k\right)  !\\
&  =\sum_{k=0}^{n}\left(  -1\right)  ^{k}\underbrace{\dbinom{n}{k}\left(
n-k\right)  !}_{\substack{=\dfrac{n!}{k!}\\\text{(by (\ref{eq.binom.fac-form}%
))}}}=\sum_{k=0}^{n}\left(  -1\right)  ^{k}\dfrac{n!}{k!}=n!\cdot\sum
_{k=0}^{n}\dfrac{\left(  -1\right)  ^{k}}{k!}.
\end{align*}


Let us summarize our results as a theorem:

\begin{theorem}
\label{thm.pie.count-der}Let $n\in\mathbb{N}$. Then, the \# of derangements of
$\left[  n\right]  $ is
\[
D_{n}=\sum_{k=0}^{n}\left(  -1\right)  ^{k}\dbinom{n}{k}\left(  n-k\right)
!=n!\cdot\sum_{k=0}^{n}\dfrac{\left(  -1\right)  ^{k}}{k!}.
\]

\end{theorem}

\begin{remark}
The sum on the right hand side is a partial sum of the well-known infinite
series $\sum\limits_{k=0}^{\infty}\dfrac{\left(  -1\right)  ^{k}}{k!}=e^{-1}$
(where $e=2.718\ldots$ is Euler's number). This is quite helpful in
approximating $D_{n}$; indeed, it is easy to see (using some simple estimates)
that%
\[
D_{n}=\operatorname*{round}\dfrac{n!}{e}\ \ \ \ \ \ \ \ \ \ \text{for each
}n>0,
\]
where $\operatorname*{round}x$ means the result of rounding a real number $x$
to the nearest integer (fortunately, since $e$ is irrational, we never get a tie).
\end{remark}

See Exercise \ref{exe.pie.dera.recs} and \cite[Chapter 1]{Wildon19} for more
about derangements.

\bigskip

\textbf{Example 3.} Like many other things in these lectures, the following
elementary number-theoretical result is due to Euler:

\begin{theorem}
\label{thm.pie.euler-tot}Let $c$ be a positive integer with prime
factorization $p_{1}^{a_{1}}p_{2}^{a_{2}}\cdots p_{n}^{a_{n}}$, where
$p_{1},p_{2},\ldots,p_{n}$ are distinct primes, and where $a_{1},a_{2}%
,\ldots,a_{n}$ are positive integers. Then,%
\[
\left(  \text{\# of all }u\in\left[  c\right]  \text{ that are coprime to
}c\right)  =c\cdot\prod_{i=1}^{n}\left(  1-\dfrac{1}{p_{i}}\right)
=\prod_{i=1}^{n}\left(  p_{i}^{a_{i}}-p_{i}^{a_{i}-1}\right)  .
\]

\end{theorem}

Note that the \# of all $u\in\left[  c\right]  $ that are coprime to $c$ is
usually denoted by $\phi\left(  c\right)  $ in number theory, and the map
$\phi:\left\{  1,2,3,\ldots\right\}  \rightarrow\mathbb{N}$ that sends each
$c$ to $\phi\left(  c\right)  $ is called \emph{Euler's totient function}.

Theorem \ref{thm.pie.euler-tot} can be proved in many ways (and a proof can be
found in almost any text on elementary number theory). Probably the most
transparent proof relies on the PIE:

\begin{proof}
[Proof of Theorem \ref{thm.pie.euler-tot} (sketched).](See \cite[proof of
Theorem 2.9.19]{19fco} for the details of this argument.) Let $U=\left[
c\right]  $. A number $u\in U$ is coprime to $c$ if and only if it is not
divisible by any of the prime factors $p_{1},p_{2},\ldots,p_{n}$ of $c$.
Again, this means that $u$ breaks all $n$ rules $1,2,\ldots,n$, where rule $i$
says \textquotedblleft thou shalt be divisible by $p_{i}$\textquotedblright.
Thus, by the \textquotedblleft rule-breaking\textquotedblright\ interpretation
of Theorem \ref{thm.pie.1}, we obtain%
\begin{align*}
&  \left(  \text{\# of all }u\in\left[  c\right]  \text{ that are coprime to
}c\right) \\
&  =\sum_{I\subseteq\left[  n\right]  }\left(  -1\right)  ^{\left\vert
I\right\vert }\underbrace{\left(  \text{\# of all }u\in\left[  c\right]
\text{ that are divisible by all }p_{i}\text{ with }i\in I\right)
}_{\substack{=\dfrac{c}{\prod_{i\in I}p_{i}}\\\text{(this is not hard to
prove)}}}\\
&  =\sum_{I\subseteq\left[  n\right]  }\left(  -1\right)  ^{\left\vert
I\right\vert }\dfrac{c}{\prod_{i\in I}p_{i}}=c\cdot\underbrace{\sum
_{I\subseteq\left[  n\right]  }\left(  -1\right)  ^{\left\vert I\right\vert
}\prod_{i\in I}\dfrac{1}{p_{i}}}_{\substack{=\prod_{i=1}^{n}\left(
1-\dfrac{1}{p_{i}}\right)  \\\text{(an easy consequence of
(\ref{eq.lem.prodrule.sum-ai-plus-bi.eq}))}}}\\
&  =\underbrace{c}_{\substack{=p_{1}^{a_{1}}p_{2}^{a_{2}}\cdots p_{n}^{a_{n}%
}\\=\prod_{i=1}^{n}p_{i}^{a_{i}}}}\cdot\prod_{i=1}^{n}\left(  1-\dfrac
{1}{p_{i}}\right)  =\left(  \prod_{i=1}^{n}p_{i}^{a_{i}}\right)  \cdot
\prod_{i=1}^{n}\left(  1-\dfrac{1}{p_{i}}\right) \\
&  =\prod_{i=1}^{n}\underbrace{\left(  p_{i}^{a_{i}}\left(  1-\dfrac{1}{p_{i}%
}\right)  \right)  }_{=p_{i}^{a_{i}}-p_{i}^{a_{i}-1}}=\prod_{i=1}^{n}\left(
p_{i}^{a_{i}}-p_{i}^{a_{i}-1}\right)  .
\end{align*}
This proves Theorem \ref{thm.pie.euler-tot}.
\end{proof}

\bigskip

\textbf{Example 4.} (This one is taken from \cite[Theorem 2.3.3]{Sagan19}.)
Recall Theorem \ref{thm.pars.odd-dist-equal}, which states that each
$n\in\mathbb{N}$ satisfies%
\[
p_{\operatorname*{odd}}\left(  n\right)  =p_{\operatorname*{dist}}\left(
n\right)  ,
\]
where%
\begin{align*}
p_{\operatorname*{odd}}\left(  n\right)   &  :=\left(  \text{\# of partitions
of }n\text{ into odd parts}\right)  \ \ \ \ \ \ \ \ \ \ \text{and}\\
p_{\operatorname*{dist}}\left(  n\right)   &  :=\left(  \text{\# of partitions
of }n\text{ into distinct parts}\right)  .
\end{align*}


We have already proved this twice, but let us prove it again.

\begin{proof}
[Third proof of Theorem \ref{thm.pars.odd-dist-equal} (sketched).]Let
$n\in\mathbb{N}$. We set $U:=\left\{  \text{partitions of }n\right\}  $.

In this proof, the word \textquotedblleft partition\textquotedblright\ will
always mean \textquotedblleft partition of $n$\textquotedblright. Thus, a
partition cannot contain any of the entries $n+1,n+2,n+3,\ldots$ (because any
of these entries would cause the partition to have size $>n$).

We want to frame the partitions of $n$ into distinct parts as rule-breakers.
We observe that%
\begin{align*}
&  \left\{  \text{partitions of }n\text{ into distinct parts}\right\} \\
&  =\left\{  \text{partitions that contain none of the entries }%
1,2,3,\ldots\text{ twice}\right\} \\
&  \ \ \ \ \ \ \ \ \ \ \ \ \ \ \ \ \ \ \ \ \left(
\begin{array}
[c]{c}%
\text{here and in the following, the word \textquotedblleft
twice\textquotedblright}\\
\text{means \textquotedblleft at least twice\textquotedblright}%
\end{array}
\right) \\
&  =\left\{  \text{partitions that contain none of the entries }%
1,2,\ldots,n\text{ twice}\right\}
\end{align*}
(since a partition cannot contain any of the entries $n+1,n+2,n+3,\ldots$). In
other words, the partitions of $n$ into distinct parts are precisely the
partitions that break all rules $1,2,\ldots,n$, where rule $i$ says
\textquotedblleft thou shalt contain the entry $i$ twice\textquotedblright%
\footnote{Once again, \textquotedblleft twice\textquotedblright\ means
\textquotedblleft at least twice\textquotedblright.}.

Thus, applying the PIE (specifically, the \textquotedblleft
rule-breaking\textquotedblright\ interpretation of Theorem \ref{thm.pie.1}),
we obtain%
\begin{align}
&  p_{\operatorname*{dist}}\left(  n\right) \nonumber\\
&  =\sum_{I\subseteq\left[  n\right]  }\left(  -1\right)  ^{\left\vert
I\right\vert }\left(  \text{\# of partitions that satisfy all rules in
}I\right) \nonumber\\
&  =\sum_{I\subseteq\left[  n\right]  }\left(  -1\right)  ^{\left\vert
I\right\vert }\left(  \text{\# of partitions that contain each of the entries
}i\in I\text{ twice}\right)  . \label{pf.thm.pars.odd-dist-equal.3rd.3}%
\end{align}


We can play the same game with $p_{\operatorname*{odd}}\left(  n\right)  $.
This time, rule $i$ says \textquotedblleft thou shalt contain the entry
$2i$\textquotedblright. Thus, again applying the PIE, we obtain%
\begin{equation}
p_{\operatorname*{odd}}\left(  n\right)  =\sum_{I\subseteq\left[  n\right]
}\left(  -1\right)  ^{\left\vert I\right\vert }\left(  \text{\# of partitions
that contain the entry }2i\text{ for each }i\in I\right)  .
\label{pf.thm.pars.odd-dist-equal.3rd.4}%
\end{equation}


Now, comparing these two equalities, we see that in order to prove that
$p_{\operatorname*{dist}}\left(  n\right)  =p_{\operatorname*{odd}}\left(
n\right)  $, it will suffice to show that%
\begin{align*}
&  \left(  \text{\# of partitions that contain each of the entries }i\in
I\text{ twice}\right) \\
&  =\left(  \text{\# of partitions that contain the entry }2i\text{ for each
}i\in I\right)
\end{align*}
for any subset $I$ of $\left[  n\right]  $.

So let $I$ be a subset of $\left[  n\right]  $. We are looking for a
bijection
\begin{align*}
&  \text{from }\left\{  \text{partitions that contain each of the entries
}i\in I\text{ twice}\right\} \\
&  \text{to }\left\{  \text{partitions that contain the entry }2i\text{ for
each }i\in I\right\}  .
\end{align*}
Such a bijection can be obtained as follows: For each $i\in I$ (from highest
to lowest\footnote{Actually, a bit of thought reveals that the order in which
we go through the $i\in I$ does not affect the result; this becomes
particularly clear if we identify each partition with the multiset of its
entries. Thus, me saying \textquotedblleft from highest to
lowest\textquotedblright\ is unnecessary.}), we remove two copies of $i$ from
the partition, and insert a $2i$ into the partition in their stead. For
example, if $I=\left\{  2,4,5\right\}  $ and $n=33$, then our bijection sends
the partition
\[
\left(  5,5,4,4,3,3,2,2,2,2,1\right)  \text{ to }\left(
10,8,4,3,3,2,2,1\right)  .
\]
(Note that the $4$ in the resulting partition is not one of the original two
$4$s, but rather a new $4$ that was inserted when we removed two copies of
$2$. On the other hand, the two $2$s in the resulting partition are inherited
from the original partition, because (unlike the bijection $A$ in our Second
proof of Theorem \ref{thm.pars.odd-dist-equal} above) our bijection only
removes two copies of each $i\in I$.)

It is easy to see that this purported bijection really is a
bijection\footnote{Its inverse, of course, does what you would expect: For
each $i\in I$, we remove a $2i$ from the partition, and insert two copies of
$i$ in its stead.}. Thus, we have found our bijection. The bijection principle
therefore yields%
\begin{align*}
&  \left(  \text{\# of partitions that contain each of the entries }i\in
I\text{ twice}\right) \\
&  =\left(  \text{\# of partitions that contain the entry }2i\text{ for each
}i\in I\right)  .
\end{align*}


We have proved this equality for all $I\subseteq\left[  n\right]  $. Hence,
the right hand sides of the equalities (\ref{pf.thm.pars.odd-dist-equal.3rd.3}%
) and (\ref{pf.thm.pars.odd-dist-equal.3rd.4}) are equal. Thus, their left
hand sides are equal as well. In other words, $p_{\operatorname*{dist}}\left(
n\right)  =p_{\operatorname*{odd}}\left(  n\right)  $. This proves Theorem
\ref{thm.pars.odd-dist-equal} again.
\end{proof}

\begin{remark}
It is worth contrasting the above four examples (in which we applied the PIE
to solve a counting problem, obtaining an alternating sum as a result) with
our arguments in Section \ref{sec.sign.intro} (in which we computed
alternating sums using sign-reversing involutions). Sign-reversing involutions
help turn alternating sums into combinatorial problems, while the PIE moves us
in the opposite direction. The two techniques are thus, in some way, inverse
to each other. This will become less mysterious once we prove the PIE itself
using a sign-reversing involution. The PIE can also be used backwards, to turn
an alternating sum into a counting problem, which is how we proved Corollary
\ref{cor.pie.count-sur.cors} above.
\end{remark}

\subsubsection{\label{subsec.sign.pie.wt}The weighted version}

The main rule of algebra is to never turn down nature's gifts. The PIE (in the
shape of Theorem \ref{thm.pie.1}) can be generalized with zero effort, so let
us do it (\cite[Theorem 7.8.9]{20f}):

\begin{theorem}
[weighted version of the PIE]\label{thm.pie.2}Let $n\in\mathbb{N}$, and let
$U$ be a finite set. Let $A_{1},A_{2},\ldots,A_{n}$ be $n$ subsets of $U$. Let
$A$ be any additive abelian group (such as $\mathbb{R}$, or any vector space,
or any ring). Let $w:U\rightarrow A$ be any map (i.e., let $w\left(  u\right)
$ be an element of $A$ for each $u\in U$). Then,%
\begin{equation}
\sum_{\substack{u\in U;\\u\notin A_{i}\text{ for all }i\in\left[  n\right]
}}w\left(  u\right)  =\sum_{I\subseteq\left[  n\right]  }\left(  -1\right)
^{\left\vert I\right\vert }\sum_{\substack{u\in U;\\u\in A_{i}\text{ for all
}i\in I}}w\left(  u\right)  . \label{eq.thm.pie.2.eq}%
\end{equation}

\end{theorem}

We can think of each value $w\left(  u\right)  $ in Theorem \ref{thm.pie.2} as
a kind of \textquotedblleft weight\textquotedblright\ of the respective
element $u$. Thus, the left hand side of the equality (\ref{eq.thm.pie.2.eq})
is the total weight of all \textquotedblleft rule-breaking\textquotedblright%
\ $u\in U$ (that is, of all $u\in U$ that satisfy $\left(  u\notin A_{i}\text{
for all }i\in\left[  n\right]  \right)  $), whereas the inner sum
$\sum_{\substack{u\in U;\\u\in A_{i}\text{ for all }i\in I}}w\left(  u\right)
$ on the right hand side is the total weight of all $u\in U$ that satisfy
$\left(  u\in A_{i}\text{ for all }i\in I\right)  $. This is why we call
Theorem \ref{thm.pie.2} the \emph{weighted version of the PIE} (or just the
\emph{weighted PIE}).

Theorem \ref{thm.pie.1} can be obtained from Theorem \ref{thm.pie.2} by taking
$w$ to be constantly $1$ (that is, by setting $w\left(  u\right)  =1$ for each
$u\in U$). Indeed, a sum of a bunch of $1$s equals the \# of $1$s being
summed, so that sums generalize cardinalities.

With Theorem \ref{thm.pie.2}, we can take any of our above four examples from
Subsection \ref{subsec.sign.pie.exas}, and introduce weights into them --
i.e., instead of asking for a number, we sum certain \textquotedblleft
weights\textquotedblright. Little question (cf. the homework): What do you get?

But we have no time for this now. Another generalization is calling!

